\documentclass[11pt]{article}
\usepackage[margin=1in]{geometry}
\usepackage{amsmath,amssymb,booktabs,array}
\usepackage[T1]{fontenc}
\usepackage{lmodern}
\usepackage[hidelinks]{hyperref}
\setlength{\parindent}{0pt}
\setlength{\parskip}{4pt}
\newcommand{\DRBaseMinusEdges}{5{,}062{,}632}
\newcommand{\DRBaseMinusStates}{547{,}021}
\newcommand{\DRBasePlusEdges}{10{,}296{,}917}
\newcommand{\DRBasePlusStates}{1{,}044{,}457}
\newcommand{\DRBaseSubsetBytes}{1{,}999{,}975{,}480}
\newcommand{\DRBaseSubsetPairs}{6{,}158}
\newcommand{\DRControlMinusEdges}{21{,}544{,}296}
\newcommand{\DRControlMinusStates}{2{,}140{,}569}
\newcommand{\DRControlPlusEdges}{32{,}980{,}090}
\newcommand{\DRControlPlusStates}{3{,}133{,}393}
\newcommand{\DRDepth}{8}
\newcommand{\DRDepthStates}{73}
\newcommand{\DRFirstCommitted}{1}
\newcommand{\DRHistoryStates}{78}
\newcommand{\DRLocalBaseEdges}{2}
\newcommand{\DRLocalBaseStates}{3}
\newcommand{\DRLocalMinusEdges}{54}
\newcommand{\DRLocalMinusStates}{34}
\newcommand{\DRLocalPlusEdges}{17}
\newcommand{\DRLocalPlusStates}{12}

\begin{document}
\begin{center}
{\Large\bfseries Response to the Third-Round Review}\\[6pt]
JBCB-1505, third major revision\\Carlos Ramirez Ovalle
\end{center}
\textbf{Manuscript title:} Locating resolution-dependent behavioral correspondence
in qualitative DNA-damage response models.

We thank the reviewer for accepting the mathematical development, which
remains unchanged. This revision focuses on the biological demonstration:
the death-receptor comparison is now the principal case and GIM a secondary
illustration. The answers below state the models, results and limitations
without requiring local documentation.

\section*{1. The GIM result is largely predetermined by the labels}

We agree with the substance of this concern. Interface C exposes mechanistic
differences already specified in the models; the method does not discover
those terminal mechanisms. Section 5 is now entitled ``GIM illustration:
interface-conditioned behavioral correspondence.'' It reports weak
correspondence at A/B and its first loss at C among the three tested
interfaces, with unchanged PN-GDDA similarity 0.9976. The abstract,
Introduction, Discussion and Conclusion no longer present GIM as the principal
scientific demonstration or as locating a unique biological resolution.
The full Petri-net diagrams and supporting evidence remain in the article.

\section*{2. Models should agree on ordinary predictions but differ in future choices}

We added a fixed-interface comparison derived from the public original model
of Calzone et al., \emph{Mathematical Modelling of Cell-Fate Decision in
Response to Death Receptor Engagement}, PLoS Computational Biology 6:e1000702
(2010), DOI: \href{https://doi.org/10.1371/journal.pcbi.1000702}{10.1371/journal.pcbi.1000702}.
Its original publication and Figures S2/S3 explicitly study removal of
CASP3-to-CASP8 feedback and the effect of transient versus sustained ligand.
DR-FB+ retains the published rules. DR-FB- removes only the CASP3 contribution
to CASP8 activation; CASP3 itself remains dynamic. The shared initial condition
has TNF=FADD=ATP=cIAP=1 and all other nodes zero. Updates are unitary
asynchronous. NFkB, CASP3 and MPT activation/deactivation are the fixed
observable events, and withdrawal changes only TNF from one to zero with
identical availability in both models. No interface was retuned to obtain
a desired relation.

\textbf{Ordinary baseline.} Complete reachable-state enumeration gives exactly
the same three stable-fate predictions under sustained TNF: survival,
apoptosis and necrosis. However, global sustained trace equality is now
refuted, not merely unresolved. An expanded search found the 12-action word
\[
\begin{split}
&\mathrm{MPT}\!\uparrow,\ \mathrm{CASP3}\!\uparrow,\
\mathrm{CASP3}\!\downarrow,\ \mathrm{NFkB}\!\uparrow,\\
&\mathrm{NFkB}\!\downarrow,\ \mathrm{MPT}\!\downarrow,\
(\mathrm{CASP3}\!\uparrow,\ \mathrm{CASP3}\!\downarrow)^3,
\end{split}
\]
where the superscript repeats the pair three times. It occurs in DR-FB+
but not DR-FB-. Independent sparse reachability and mCRL2 confirm this
decision; a 53-update accepting path was also checked against the full
28-node rules. Consequently DR-FB- cannot weakly simulate DR-FB+, and the
global baseline systems are not weakly bisimilar. The reverse trace inclusion
remains undecided. These findings are stated in the abstract, Figure 1, main
Results and Discussion. The word is a model prediction, not experimental evidence.

\textbf{Specific shared history.} The automatic search identifies a jointly
reachable retained state after eight updates: TNFR, DISC\_TNF, CASP8, BAX,
MOMP, Cyt\_c, apoptosome and CASP3 activation. Only the last update is visible
at the declared interface, so the common observable history is
$h=\texttt{CASP3\_up}$. At this state, TNF, FADD, ATP, cIAP and all eight
listed components are one; the other retained components are zero.
The three omitted nodes, NonACD, Apoptosis and Survival, are unobserved
output-only copies with no feedback. Their deletion preserves weak behavior,
and the full 28-node conditioned systems were also checked directly.

\textbf{Exact conditioned baseline and intervention result.} From the
specified shared state, the two sustained continuations each have three
states and two silent edges. Their exact observable language is
$\{\epsilon\}$, and they are strongly and weakly bisimilar. Once withdrawal
is permitted, the continuations have \DRLocalPlusStates{}/\DRLocalPlusEdges{}
and \DRLocalMinusStates{}/\DRLocalMinusEdges{} states/edges. Writing them
$K_+$ and $K_-$, with $X\preceq_wY$ meaning that $Y$ simulates $X$, we find
\[
K_+\preceq_wK_-,\qquad K_-\not\preceq_wK_+,\qquad K_+\not\sim_wK_-.
\]
The word \texttt{withdraw\_TNF, CASP3\_down} occurs only in $K_-$. The
independent game certificate eliminates the initial pair only after covering
all possible defending weak replies. It does not infer a failed global
relation from one arbitrarily chosen successor pair. Production predicates,
the independent game and mCRL2 equivalence checks agree.

After immediate withdrawal from this shared state, $K_+$ retains CASP3 on
every continuation and has only an apoptotic terminal fate. $K_-$ has only
the naive terminal signature. Both continuation graphs are acyclic; their
post-withdrawal outcome does not require an added fairness assumption.
``Naive'' is a logical marker signature, not evidence that a cell can recover
after executing apoptosis.

\textbf{Visible history versus hidden state.} The history $h$ is compatible
with 78 internal states in each model, not just the selected state. Complete
aggregation after immediate withdrawal yields
\[
\begin{split}
\mathrm{Futures}_+(h)&=\{\mathrm{apoptosis},\mathrm{naive},\mathrm{necrosis}\},\\
\mathrm{Futures}_-(h)&=\{\mathrm{naive},\mathrm{necrosis}\}.
\end{split}
\]
Thus a future-set difference also exists at the shared-history level, but
CASP3 activation alone does not establish unique commitment in every compatible
feedback-positive state. Section 4 distinguishes this aggregate claim from
the fully specified state's persistent-CASP3 certificate.

\textbf{What is not established.} The complete intervention-aware languages
are unequal: exact counterexamples refute both global trace inclusions and,
consequently, both global weak simulations. For example,
\texttt{CASP3\_up, withdraw\_TNF, CASP3\_down} is possible only in DR-FB-.
In the other direction, \texttt{withdraw\_TNF, CASP3\_up, CASP3\_down,
CASP3\_up} is possible only in DR-FB+.
Trace comparison can therefore distinguish the models both before and after
withdrawal is included. We do not present this case as proving that branching alone
distinguishes globally trace-equivalent biological models. It supplies the
weaker, biologically grounded result of matching baseline endpoints and exact
conditioned sustained continuations, followed by different intervention
futures. The stronger biological counterpart of the accepted early/late-choice
example remains an open limitation.

\section*{3. Identify a meaningful perturbation or observation}

The new Figure 1 and Section 4 identify the logic of a pulse/withdrawal
comparison at a CASP3-positive stage. Subsequent joint readout of CASP3,
NFkB and MPT can distinguish persistent apoptotic activity from deactivation
or movement toward the necrotic branch; ATP supports terminal classification.
Monitoring all three matters because a path with MPT activation before CASP3
loss is not the same as isolated CASP3 deactivation after withdrawal.
The history alone does not identify the internal witness, so the state-specific
prediction would require additional state assessment. We provide neither
unsupported physical timings nor an invented selective inhibitor protocol.

The source publication already studied ligand withdrawal and the feedback
mechanism. Our contribution is the reproducible shared-state/history
construction, future-set comparison and checked relational explanation.
This is a recovery and characterization of a published phenomenon, not a new
wet-lab discovery. No new experimental confirmation is claimed.

\section*{4. Qualify the claimed biological boundary}

The GIM interpretation is now centrally restricted to the first loss of
correspondence among Interfaces A, B and C. It is not a unique, globally
minimal or experimentally established biological boundary. No exhaustive
search over observational interfaces was performed. Section 8 separates this
interface-change question from the fixed-interface withdrawal question.

\section*{5. Distinguish specified knowledge from computed findings}

The main article now includes an explicit evidence-status table:
the feedback mechanism is a published assumption; the marker interface and
withdrawal semantics are declared inputs; stable-fate sets, relation
decisions and witness histories are computed outputs; the discriminating
observation is model-derived; wet-lab confirmation is absent. The local
prospective configuration is not presented as externally registered.
Resource-driven changes remove only provably unobserved output sinks and use
compact storage; they do not change observables or biological rules.

\section*{6. Preserve the accepted mathematics and ensure reproducibility}

All accepted definitions, propositions, proofs, the fixed-point algorithm and
the early/late-choice example are retained, with an automated source comparison
checking their preservation. The new witness engine is tested on the accepted
example, including rejection of a certificate that omits a defending reply.
The new biological computation independently verifies compact execution,
terminal reachability, marker persistence and both simulation directions.
The eight-update commitment state is found by an exhaustive depth scan; at
that depth, one of 73 states has persistent CASP3 and a unique apoptotic
terminal fate after withdrawal. Counts are not biological probabilities,
and update depth is not physical time.

Section 7 and Supplementary Validation report the computational checks,
source provenance, semantics, limits and certificates.
\end{document}
